\section{Bell Inequalities Revisited: Structural vs Dynamical Assumptions}
  \label{sec:abstract-math-model}

  Bell’s theorem is commonly presented as a no-go result for local realist theories.
  More precisely, it establishes that no model satisfying a specific set of probabilistic
  assumptions can reproduce all quantum correlations.
  It is therefore essential to distinguish carefully between the assumptions entering
  the derivation of Bell inequalities and the physical principles they are often taken to represent.

  The derivation of Bell inequalities relies on three key assumptions.
  First, outcome probabilities are assumed to be well-defined and conditioned on a
  complete specification of relevant variables.
  Second, joint probabilities are assumed to factorize for spacelike separated
  measurements, reflecting the absence of direct causal influence between distant measurement settings.
  Third, the conditioning variables are assumed to be statistically independent of the measurement choices.
  Violations of Bell inequalities imply that at least one of these assumptions fails.

  In many discussions, the failure of Bell inequalities is interpreted as evidence for a breakdown of
  locality~\cite{Bell1964}.
  However, locality in Bell’s sense is a condition on the factorization of effective
  probabilities, not a direct statement about the causal structure of the underlying
  theory.
  Bell inequalities constrain models in which observable outcomes can be represented
  as functions of local settings and a shared set of variables that fully determine
  those outcomes.
  They do not directly constrain the existence or absence of superluminal interactions.

  Within the framework considered here, the failure of Bell inequalities originates
  from the breakdown of effective factorization.
  As shown in the previous section, non-injective mappings between underlying
  configurations and observable outcomes generically prevent the construction of a
  complete set of outcome-defining variables at the effective level.
  Although the underlying description may be local and deterministic, the observable
  description does not admit a factorized probabilistic representation.

  This distinction clarifies the sense in which Bell’s theorem applies.
  Bell inequalities are valid only for effective descriptions in which observable
  outcomes can be parameterized by a set of variables that renders the projection from
  underlying configurations injective or equivalently factorizable.
  When this condition is not met, the assumptions required for the derivation of Bell
  inequalities are violated independently of any consideration of causal nonlocality.

  From this perspective, violations of Bell inequalities do not compel the introduction
  of nonlocal dynamics or retrocausal mechanisms.
  They instead signal a structural limitation of effective descriptions in which
  observable outcomes are treated as independently specifiable local properties.
  Nonlocal correlations arise because joint outcomes must be globally consistent with a
  single non-factorizable effective description, not because information propagates
  between distant regions.

  The present analysis therefore reframes Bell’s theorem as a constraint on a class of
  descriptive models rather than as a direct statement about the ontology or dynamics of physical reality.
  Bell inequalities test the compatibility of observed correlations with injective and
  factorizable effective descriptions.
  When non-injective projections are involved, Bell-type correlations may arise while
  preserving locality at the underlying level.
